% !TEX program = pdflatex
\documentclass[11pt]{article}

\usepackage[a4paper,margin=1in]{geometry}
\usepackage{amsmath,amssymb,amsthm,mathtools}
\usepackage{booktabs}
\usepackage{microtype}
\usepackage[T1]{fontenc}
\usepackage{lmodern}
\usepackage{graphicx}
\usepackage{flafter}
\usepackage{listings}
\usepackage{xcolor}
\usepackage{hyperref}

\graphicspath{{figures/}}

\hypersetup{colorlinks=true,linkcolor=blue,citecolor=blue,urlcolor=blue}

\definecolor{codebg}{RGB}{255,250,225}  % soft yellow behind code

\lstset{
  basicstyle=\ttfamily\small,
  backgroundcolor=\color{codebg},
  keywordstyle=\color{blue!60!black},
  commentstyle=\color{black!55},
  stringstyle=\color{green!35!black},
  language=Python,
  showstringspaces=false,
  columns=fullflexible,
  keepspaces=true,
  frame=single,
  rulecolor=\color{black!25},
  xleftmargin=0pt,
  framexleftmargin=4pt,
  framexrightmargin=4pt,
}

% The numbers of the opening example and of the real data, written by python/make_figures.py.
\input{figures/example.tex}
\input{figures/scd14.tex}
\input{figures/classo.tex}

\newtheorem{proposition}{Proposition}
\theoremstyle{remark}
\newtheorem{remark}{Remark}

\newcommand{\trans}{^{\top}}
\newcommand{\Sig}{\Sigma}
\newcommand{\R}{\mathbb{R}}
\newcommand{\one}{\mathbf{1}}
\newcommand{\F}{\mathcal{F}}
\newcommand{\B}{\mathcal{B}}

\title{The Constrained LASSO by the Critical Line Algorithm\\[6pt]
{\normalsize Two ways to hand a regression to a portfolio optimiser}}
\author{Thomas Schmelzer\\[2pt]
{\small Jebel Quant Research, Abu Dhabi}\\[1pt]
{\small ORCID \href{https://orcid.org/0009-0009-9771-3501}{0009-0009-9771-3501}}}
\date{18 October 2026}

\begin{document}

\maketitle

\begin{abstract}
The LASSO was introduced, and is still mostly solved, with no constraint other than its
$\ell_1$ budget. Many regressions need more. In a regression on compositional data the
coefficients must sum to zero, and \texttt{lars\_path} in scikit-learn has no way to say
so. The usual workaround, regressing on log-ratios against a reference part, solves a
different problem for each choice of reference: on a simulated example only
\exAnywhere{} of \exParts{} choices select the right variables at any $\lambda$, and on
real microbiome data at most \rdMaxAnywhere{} of \rdGenera{} choices ever select what
the zero-sum LASSO selects. A short proposition explains why: the two paths meet only where
the reference part is unselected and its generalised correlation is zero, at isolated
points. This note computes the
exact path of the constrained LASSO with Markowitz's Critical Line Algorithm (CLA),
unmodified. That works because a LASSO is a portfolio problem, as Schmelzer and Hastie
showed, and the CLA has handled linear constraints since 1956. There are two ways to
pose the regression as a portfolio. In the \emph{penalty form} every unit of gross
position costs $\lambda$ units of return, and the CLA's parameter is $\lambda$ itself;
bounds and linear equality constraints carry over. In the \emph{budget form} the $\ell_1$
ball becomes Markowitz's long-only, fully invested budget with a cash asset; constraints
such as $\one\trans\beta=0$ carry over, and the path ends when cash enters. On the
diabetes data, a $p>n$ design and compositional data, the CLA reproduces every LARS
knot to $10^{-13}$ and meets the LASSO's optimality conditions to $10^{-15}$.

\medskip\noindent
\textbf{Keywords:} constrained LASSO; critical line algorithm; solution path; least angle
regression; compositional data; portfolio optimisation; reproducibility
\end{abstract}

\section{A regression \texttt{lars\_path} cannot fit}
\label{sec:example}

Tibshirani \cite{tibshirani1996} introduced the LASSO as least squares with a single
constraint, a budget on the $\ell_1$ norm of the coefficients:
\[
\min_\beta\ \tfrac12\|y-X\beta\|_2^2\quad\text{s.t.}\quad\|\beta\|_1\le t,
\qquad\text{or equivalently}\qquad
\min_\beta\ \tfrac12\|y-X\beta\|_2^2+\lambda\|\beta\|_1 .
\]
LARS \cite{efron2004} and the homotopy method \cite{osborne2000} compute the solution
for every $t$, or every $\lambda$, at once. The solution is piecewise linear in the
parameter, so the whole \emph{path} is a finite list of \emph{knots}. scikit-learn's
\texttt{lars\_path} \cite{sklearn} takes $X$, $y$ and one option, \texttt{positive},
for nonnegative coefficients. There is no place for any other constraint.

Many regressions need one. Here is an example in the style of Lin et al.\
\cite{lin2014}. Each of $n=60$ samples is a composition of $p=\exParts$ parts, such as
the relative abundances of bacteria in a gut sample, and $z_{ij}$ is the log of the
share of part $j$ in sample $i$. The model is
\[
y_i=\sum_j\beta_j z_{ij}+\varepsilon_i,\qquad\sum_j\beta_j=0 .
\]
The constraint is what makes the model meaningful. A composition is only known up to
scale: measuring a sample in other units multiplies all its parts by some $c$, which
adds $\log c$ to every $z_{ij}$ and $\log c\sum_j\beta_j$ to the fit. Only when the
coefficients sum to zero does the fit not depend on the units. In the simulation five
parts have nonzero true coefficients, and the goal is to find them.

\texttt{lars\_path} offers two ways out, and both fail. With the log-compositions in
\texttt{Z}:
\begin{lstlisting}
alphas, _, coefs = lars_path(Z, y, method="lasso")   # ignores sum(beta) = 0
others = np.delete(np.arange(p), r)                  # or: drop a reference part r
alphas, _, coefs = lars_path(Z[:, others] - Z[:, [r]], y, method="lasso")
\end{lstlisting}
\begin{itemize}\setlength{\itemsep}{0pt}
\item \emph{Ignore the constraint.} At $\lambda=\exLambda$ the coefficients
\texttt{lars\_path} returns sum to $\exPlainSum$, not zero, so the fitted model changes
with the units of the data.
\item \emph{Eliminate a reference part.} Pick a part $r$, substitute
$\beta_r=-\sum_{j\ne r}\beta_j$, and regress on the log-ratios $z_j-z_r$. That
regression is unconstrained, so \texttt{lars\_path} runs. But its penalty is
$\sum_{j\ne r}|\beta_j|$: the reference part is not penalised, so each choice of $r$ is
a different problem with a different answer. Figure~\ref{fig:example}, right, shows all
\exParts{} answers at $\lambda=\exLambda$. Only \exMatches{} of them select the five
parts of the true model, and the coefficients differ from those of the zero-sum LASSO
by up to \exWorst\% of the largest one. Since the reference goes unpenalised, the same
$\lambda$ does not shrink the two problems alike, so we also ask the weaker question
whether a log-ratio path selects the five parts at \emph{any} $\lambda$. Only
\exAnywhere{} of the \exParts{} do. Section~\ref{sec:reference} proves that this is no
accident of the data: a log-ratio path meets the zero-sum path only at isolated points.
\end{itemize}
A third option is correct but blind: solve the constrained problem with a general
convex solver on a grid of $\lambda$. That gives points, not the path, and it cannot
tell where a variable enters between two grid points.

On real data the workaround fares worse still. In the microbiome data of
Section~\ref{sec:real}, with \rdGenera{} genera, take any of the first \rdSegments{}
segments of the zero-sum path: at most \rdMaxAnywhere{} of the \rdGenera{} possible
references select the same genera at any $\lambda$ of their own paths.

None of this is news to people who fit such models. The zero-sum constraint of Lin et
al.\ \cite{lin2014} treats all parts alike precisely so that no part has to serve as the
reference. Altenbuchinger et al.\ \cite{zerosum} show that a poor choice of reference
degrades the LASSO, and their \texttt{zeroSum} package fits zero-sum elastic nets by
coordinate descent, again on a grid of $\lambda$. The Python package \texttt{c-lasso}
\cite{classo} is built for exactly this model, the LASSO under $C\beta=0$ with
log-contrast regression as its leading example, and includes a path algorithm that
reports the solution on a grid of $\lambda$; Section~\ref{sec:evidence} compares it with
the exact path. The point of the example is narrower:
the standard path software stops at Tibshirani's problem, so the exact path of a
constrained LASSO needs a different algorithm.

\begin{figure}[t]
\centering
\includegraphics[width=\textwidth]{example.pdf}
\caption{The opening example. Left: the exact path of the zero-sum LASSO, computed by
the CLA in the budget form (lines), with an interior-point solve at the midpoint of every
segment (circles). The five parts of the true model enter first. Right: the coefficients
at $\lambda=\exLambda$ (dashed line on the left). Dots: the zero-sum LASSO. Dashes:
\texttt{lars\_path} on log-ratios, one answer for each of the \exParts{} possible
reference parts. The ticks mark the parts in the true model.}
\label{fig:example}
\end{figure}

The left panel of Figure~\ref{fig:example} is the exact path, knot by knot. It was
computed by an algorithm that knows nothing about regression: the Critical Line
Algorithm (CLA) that Markowitz \cite{markowitz1956} wrote for portfolios. The same path
is one call in the open-source package \texttt{cvxcla} \cite{cvxcla}, which traces
homogeneous equality constraints with the rescaling of Section~\ref{sec:budget}:
\begin{lstlisting}
from cvxcla import Lasso
knots = Lasso(Z, y, a=np.ones((1, p))).path   # each knot: lam, beta, active
\end{lstlisting}

\paragraph{Why a portfolio algorithm} Schmelzer and Hastie \cite{perspective} prove that
the CLA and the constrained LASSO \cite{james2020,gaines2018} trace the same curve.
Read a regression as a portfolio problem, as in Table~\ref{tab:dictionary}: the
coefficients are positions, $X\trans X$ is the covariance of the assets, $X\trans y$ is
their expected return and $\lambda$ is a risk tolerance. Portfolio managers have always
needed constraints, such as a budget, position limits and sector exposures, so
Markowitz built them into the CLA from the start. Once a LASSO is posed as a portfolio,
its constraints come along at no cost. A companion note \cite{lassonote} used the same
identity the other way, computing efficient frontiers with a LASSO solver.

\begin{table}[t]
\centering
\small
\begin{tabular}{ll}
\toprule
regression & portfolio \\
\midrule
coefficient $\beta_j$ & position in asset $j$ \\
$X\trans X$ & covariance of the asset returns \\
$X\trans y$ & expected returns (penalty form: a benchmark) \\
$\|\beta\|_1$ & gross exposure \\
penalty $\lambda$ & risk tolerance $\gamma$ (penalty form) \\
bound $\|\beta\|_1\le t$ & fully invested budget, $\gamma=1/t$ (budget form) \\
$A\beta=b$, $l\le\beta\le u$ & budget rows and position limits \\
knot of the LASSO path & turning point of the CLA \\
\bottomrule
\end{tabular}
\caption{The dictionary between the two problems \cite{perspective}.}
\label{tab:dictionary}
\end{table}

Exact paths of the constrained LASSO are not new. Gaines, Kim and Zhou
\cite{gaines2018} derive a path algorithm for it from its optimality conditions, which
tracks the active set and the multipliers of the linear constraints as $\lambda$
decreases. The CLA does the same bookkeeping: both are active-set methods that follow
one family of quadratic programs, which is what the identity of \cite{perspective}
leads one to expect. The difference is where the algorithm comes from. Here it is not
derived for the LASSO at all. Markowitz's algorithm runs unchanged, so seventy years of
work on it, and its implementations, such as \texttt{cvxcla} \cite{cvxcla}, carry over
to the LASSO. The price is generality: the encodings below need a nondegenerate start
and leave out general inequality rows (Section~\ref{sec:limits}). The work is in the
encoding, and there are two:
\begin{enumerate}\setlength{\itemsep}{0pt}
\item the \emph{penalty form} (Section~\ref{sec:penalty}) keeps $\lambda$ as the
parameter, and takes coefficient bounds and general equality constraints $A\beta=b$;
\item the \emph{budget form} (Section~\ref{sec:budget}) keeps the $\ell_1$ budget, as
Markowitz's long-only, fully invested problem with a cash asset. It takes homogeneous
constraints such as $\one\trans\beta=0$, and it is the form that solves the example.
\end{enumerate}
Neither reduction is deep, and the note claims no new theorem. What it adds is a
working instance: each form is a few lines on top of a CLA of about a hundred lines, and
Section~\ref{sec:evidence} shows that together they reproduce LARS and an
interior-point solver to rounding error. Section~\ref{sec:limits} sets out where both
forms stop. Readers who know the CLA can skip Section~\ref{sec:cla}.

\paragraph{Notation} The design is $X\in\R^{n\times p}$, with $n$ observations and $p$
predictors, and the response is $y\in\R^n$. Write $G=X\trans X$ and $c=X\trans y$. The
constrained LASSO is
\begin{equation}
\label{eq:lasso}
\mathrm{L}_\lambda:\quad
\min_\beta\ \tfrac12\|y-X\beta\|_2^2+\lambda\|\beta\|_1
\quad\text{s.t.}\quad A\beta=b,\quad l\le\beta\le u,
\end{equation}
with $A\in\R^{m\times p}$ and $l_j\le0\le u_j$, infinite entries allowed. With no rows
and infinite bounds it is Tibshirani's LASSO. We assume, as \cite{perspective} does,
that $\mathrm{L}_\lambda$ has a unique minimiser $\beta(\lambda)$ for every
$\lambda>0$, which is the generic case.

\section{The Critical Line Algorithm}
\label{sec:cla}

In words, the CLA starts with the riskiest portfolio, the one with the highest expected
return, and lowers the risk tolerance step by step. Between steps the optimal
portfolio moves along a straight line. A step ends when a position hits a limit or an
asset at a limit wants to come off it. These corners are the \emph{turning points};
the LASSO calls them knots.

Precisely, the CLA solves the family of quadratic programs
\begin{equation}
\label{eq:cla}
\mathrm{P}_\gamma:\quad
\min_x\ \tfrac12\,x\trans\Sig x-(q+\gamma\mu)\trans x
\quad\text{s.t.}\quad Ax=b,\quad l\le x\le u,
\end{equation}
for every $\gamma$ from $+\infty$ down to $0$, with $\Sig\succeq0$. Here $x$ is the
vector of holdings, $\Sig$ the covariance, $\mu$ the expected returns and $\gamma$ the
risk tolerance. The term $q$ is a benchmark. When $q=\Sig\bar x$, the objective is
$\tfrac12(x-\bar x)\trans\Sig(x-\bar x)-\gamma\mu\trans x$ up to a constant, the
tracking variance against a portfolio $\bar x$ traded off against expected return. This
is the same CLA after the change of variables $x-\bar x$, whose only effect is to shift
the bounds and the right-hand side. With $q=0$ it is Markowitz's problem exactly.

\paragraph{Segments} Write
$g=\Sig x-q-\gamma\mu+A\trans\nu$ for the gradient of the Lagrangian. A point $x$ with
multipliers $\nu$ is optimal when $g_j=0$ on the \emph{free} coordinates
$\F=\{j:l_j<x_j<u_j\}$, $g_j\ge0$ where $x_j=l_j$, and $g_j\le0$ where $x_j=u_j$. Call
the coordinates at a bound \emph{blocked}, $\B$. With $\F$ and $\B$ fixed, $x_\F$ and
$\nu$ solve the bordered system
\begin{equation}
\label{eq:border}
\begin{bmatrix}\Sig_{\F\F}&A_\F\trans\\A_\F&0\end{bmatrix}
\begin{bmatrix}x_\F\\\nu\end{bmatrix}
=\begin{bmatrix}q_\F-\Sig_{\F\B}x_\B\\b-A_\B x_\B\end{bmatrix}
+\gamma\begin{bmatrix}\mu_\F\\0\end{bmatrix},
\end{equation}
so $x_\F=\alpha+\gamma\delta$ and $\nu=\nu_0+\gamma\nu_1$ are affine in $\gamma$. The
gradient on the blocked coordinates is affine too, $g_\B=g_0+\gamma g_1$. One
factorisation and two right-hand sides give the whole segment.

\paragraph{Events} Lowering $\gamma$, the segment ends at the largest $\gamma$ below the
current one at which either a free coordinate reaches a bound,
$\gamma=(l_i-\alpha_i)/\delta_i$ for $\delta_i>0$ or $(u_i-\alpha_i)/\delta_i$ for
$\delta_i<0$, or a blocked coordinate's gradient changes sign,
$\gamma=-g_{0j}/g_{1j}$ for $g_{1j}>0$ at a lower bound or $g_{1j}<0$ at an upper one.
The coordinate concerned changes sides, \eqref{eq:border} is solved again, and the walk
continues. If there is no event the last segment runs to $\gamma=0$.

\paragraph{The start} As $\gamma\to\infty$ the linear term dominates, and the path
starts at a vertex maximising $\mu\trans x$ over the feasible set. A linear program
finds it. If the maximiser is unique and has exactly $m$ coordinates strictly inside
their bounds, the bordered matrix in \eqref{eq:border} is nonsingular and the walk can
begin. This is the nondegenerate case. The implementation checks the count and refuses
to start otherwise.

\section{The penalty form}
\label{sec:penalty}

The first encoding reads the penalty $\lambda\|\beta\|_1$ as a cost: every unit of
gross position, long or short, costs $\lambda$ units of return. To make that a linear
cost, each coefficient is split into a long and a short position, both nonnegative.
The constraints on $\beta$ become constraints on the positions, unchanged in kind.

Formally, let $S$ have a column $e_j$ for every $j$ with $u_j>0$ and a column $-e_j$
for every $j$ with $l_j<0$, so $\beta=Sz$ with $z\ge0$, and let $\bar u$ collect the
matching bounds $u_j$ and $-l_j$. A coefficient whose bounds are both zero is fixed and
can be dropped, so every $j$ has at least one column and $S$ has full row rank.

\begin{proposition}[Penalty as cost]
\label{prop:penalty}
For $\lambda>0$, $\beta(\lambda)=Sz(\lambda)$, where $z(\lambda)$ solves the CLA
instance $\mathrm{P}_\gamma$ with
\[
\Sig=S\trans GS,\quad \mu=-\one,\quad q=S\trans c,\quad
\text{rows } ASz=b,\quad 0\le z\le\bar u,\quad \gamma=\lambda .
\]
Moreover $q$ lies in the range of $\Sig$, so the benchmark reading of
Section~\ref{sec:cla} applies: the program tracks a least-squares fit and pays
$\lambda$ for every unit of gross position.
\end{proposition}

\begin{proof}
Up to the constant $\tfrac12\|y\|^2$, the CLA objective at $\gamma=\lambda$ is
$\tfrac12\|y-XSz\|^2+\lambda\one\trans z$. For a fixed $\beta$ feasible in
\eqref{eq:lasso}, the smallest $\one\trans z$ over $z\ge0$ with $Sz=\beta$ and $z\le\bar u$
is $\|\beta\|_1$. It is attained by putting $\beta_j$ on the column of its sign, which
respects $\bar u$ because $l_j\le\beta_j\le u_j$. The minimum of the CLA program over
$z$ is therefore the minimum of $\mathrm{L}_\lambda$ over $\beta$. At a minimiser no
coordinate holds both signs, since lowering both by the same amount keeps $\beta$ and
reduces the cost by $2\lambda$ times that amount. Hence $Sz(\lambda)$ minimises
$\mathrm{L}_\lambda$, and by uniqueness equals $\beta(\lambda)$. For the range claim,
$c=X\trans y$ lies in the range of $G$, say $c=Gv$. Since $S$ has full row rank,
$v=S\bar x$ for some $\bar x$, and then $q=S\trans GS\bar x=\Sig\bar x$, where $S\bar x$
is a least-squares fit.
\end{proof}

In code the form is three lines on top of a function \texttt{cla} that implements
Section~\ref{sec:cla}, here with both signs allowed for every coefficient:
\begin{lstlisting}
S = np.hstack([np.eye(p), -np.eye(p)])        # beta = S z with z >= 0
run = cla(S.T @ G @ S, -np.ones(2 * p), A @ S, b, np.zeros(2 * p), cap, q=S.T @ c)
beta = run.x @ S.T                            # one row per knot, at lam = run.gamma
\end{lstlisting}

The covariance $S\trans GS$ is singular whenever both signs of a coefficient are
allowed, because the long and short columns cancel. The bordered system
\eqref{eq:border} is nonetheless nonsingular on every segment with $\lambda>0$: the
argument in the proof keeps the two halves of a coefficient out of $\F$ together, and
under the general-position assumption of \cite{perspective} the columns of $X$ on the
support are independent. The halves come apart only at $\lambda=0$, where the gradient
of a free half's twin is $2\lambda$ and vanishes. In floating point that zero can land
at $\gamma\approx10^{-16}$ instead, so the implementation treats event parameters below
$10^{-10}$ times the first one as zero.

\paragraph{The start} Without equality rows the starting vertex is $z=0$: all
positions are closed, because each costs more than it can earn. The first event is the
first entry, at $\gamma=\max_j|c_j|$ over the allowed signs, which is LARS's
$\lambda_{\max}$. With equality rows the start is the linear program $\min\one\trans z$
over $ASz=b$, $0\le z\le\bar u$, the point of smallest $\ell_1$ norm on the constraint
set. For generic $A$ and $b\ne0$ it is a unique vertex with $m$ nonzeros. The path then
holds still on its first segment, $\delta=0$, while the multipliers absorb the growing
penalty, until a second coefficient enters.

Figure~\ref{fig:path} shows the penalty form on the diabetes data \cite{efron2004},
where there are no constraints and LARS can check every knot.

\begin{figure}[t]
\centering
\includegraphics[width=0.66\textwidth]{path.pdf}
\caption{The diabetes LASSO path \cite{efron2004}, computed by the CLA in the penalty
form (lines). Circles: the knots \texttt{lars\_path} returns \cite{sklearn}. All 13
knots agree to $2\times10^{-13}$, relative.}
\label{fig:path}
\end{figure}

\section{The budget form}
\label{sec:budget}

The penalty form cannot solve the opening example: with $b=0$ its start is degenerate
(Section~\ref{sec:limits}). The second encoding starts instead from Tibshirani's
original statement, the budget $\|\beta\|_1\le t$, and divides by $t$. The split
coefficients then add up to at most one, which is Markowitz's own setting: a long-only
portfolio that invests a budget of one. A cash asset, with no risk and no return, holds
whatever part of the budget the regression does not use.

The bound form has the same path as $\mathrm{L}_\lambda$, indexed by
$t=\|\beta(\lambda)\|_1$ instead of $\lambda$ \cite{perspective}.

\begin{proposition}[Budget as simplex]
\label{prop:budget}
Suppose $b=0$ and $l_j,u_j\in\{0,\pm\infty\}$, so that every constraint is homogeneous.
Write $\beta=tSw$ with $w\ge0$, and add a riskless asset $w_0\ge0$ with
$\one\trans w+w_0=1$. Then for $t>0$ the bound-form LASSO is the CLA instance
\[
\Sig=\begin{bmatrix}S\trans GS&0\\0&0\end{bmatrix},\quad
\mu=\begin{bmatrix}S\trans c\\0\end{bmatrix},\quad q=0,\quad \gamma=1/t,
\]
with rows $\one\trans w+w_0=1$ and $ASw=0$ and bounds $w,w_0\ge0$. This is
Markowitz's long-only, fully invested problem with a cash asset of zero variance and
zero return. If $\eta$ is the multiplier of the budget row and $\nu$ those of
$ASw=0$, then $\lambda=t\eta$ and the LASSO's multipliers for $A\beta=0$ are $t\nu$.
\end{proposition}

\begin{proof}
As in Proposition~\ref{prop:penalty}, $\|\beta\|_1\le t$ is $\one\trans z\le t$ for the
split $z=tw$, and $w_0$ is its slack. The constraint $ASz=0$ is unchanged by the scale
$t$, which is where homogeneity is needed. A nonzero right-hand side would become
$b/t=\gamma b$, a constraint that moves with the parameter. Substituting $z=tw$ gives the
objective $t^2(\tfrac12w\trans S\trans GSw-\gamma c\trans Sw)$ plus a constant, and
dividing by $t^2>0$ gives $\mathrm{P}_\gamma$. For the multipliers, multiply the CLA's
stationarity row for a free coordinate,
$(S\trans GSw)_j-\gamma(S\trans c)_j+\eta+(S\trans A\trans\nu)_j=0$, by $t$. This gives
$(S\trans(G\beta-c))_j+t\eta+(S\trans A\trans t\nu)_j=0$, the LASSO's stationarity
condition with penalty $t\eta$ and constraint multipliers $t\nu$. The cash row reads
$\eta\ge0$ while cash is held at zero.
\end{proof}

In code, cash is one more asset with no variance and no return, and the run stops when
it enters:
\begin{lstlisting}
k = 2 * p + 1                                    # long, short, cash
Sig = np.zeros((k, k)); Sig[:-1, :-1] = S.T @ G @ S
mu = np.append(S.T @ c, 0.0)
rows = np.vstack([np.ones(k), np.hstack([A @ S, np.zeros((m, 1))])])
run = cla(Sig, mu, rows, np.append(1.0, np.zeros(m)),
          np.zeros(k), np.full(k, np.inf), stop=k - 1)
beta = run.x[:, :-1] @ S.T / run.gamma[:, None]  # beta = t S w with t = 1/gamma
\end{lstlisting}

The two ends of the CLA path have plain meanings. At $\gamma=\infty$, $t=0$, the start
is the linear program $\max c\trans Sw$ over the simplex: put the whole budget on the
best asset. Without equality rows that is the predictor most correlated with the
response, entering with the sign of its correlation and with
$\lambda=\lim\eta/\gamma=\max_j|c_j|$, just as in LARS. Under $\one\trans\beta=0$ it is
one long and one short position of weight $\tfrac12$ each, the most and the least
correlated predictors, entering together. That is why two coefficients leave zero at
once on the left of Figure~\ref{fig:example}. Both starts are generically unique and
nondegenerate.

At the other end, cash leaves its bound when its gradient $\eta$ reaches zero, so
$\lambda=0$ and the budget goes slack: the regression no longer wants a larger $\ell_1$
norm, and the LASSO path is over. Past that point the CLA still has a path, the ray
$w=\beta(0)/t$ diluted by cash, but the twin halves' gradients vanish identically along
it. The implementation therefore ends the run when cash enters, and lets that event win
a tie. The tie matters when $p>n$: there every correlation vanishes at $\lambda=0$
together with $\eta$, and rounding alone would decide which event the walk takes.
Figure~\ref{fig:portfolio} shows the diabetes path in these terms.

The budget form gives the same knots as the penalty form only where the map from
$\lambda$ to $t$ is one-to-one, which holds where $\|\beta(\lambda)\|_1$ is strictly
decreasing \cite[Proposition~1]{perspective}. That fails only if the support's sign
vector lies in the span of the active constraint normals. Under $\one\trans\beta=0$ that
would require every coefficient on the support to have the same sign, which is
impossible for a nonzero $\beta$ that sums to zero. So the zero-sum constraint is safe.

\medskip\noindent
\fbox{\begin{minipage}{\dimexpr\linewidth-2\fboxsep-2\fboxrule\relax}
\textbf{Algorithm 1: the constrained LASSO path from a CLA.}
\begin{enumerate}\setlength{\itemsep}{0pt}
\item Form $G=X\trans X$ and $c=X\trans y$. Build $S$ with a column $e_j$ for each $j$
with $u_j>0$ and a column $-e_j$ for each $j$ with $l_j<0$.
\item If every constraint is homogeneous, $b=0$ and $l_j,u_j\in\{0,\pm\infty\}$, use the
budget form: assets $(w,w_0)$ with covariance $\operatorname{diag}(S\trans GS,0)$ and
expected returns $(S\trans c,0)$, rows $\one\trans w+w_0=1$ and $ASw=0$, and $w,w_0\ge0$.
Run the CLA from $\gamma=\infty$ and stop when cash leaves zero. A turning point at
$\gamma$ with budget multiplier $\eta$ is the knot $\beta=Sw/\gamma$ at
$\lambda=\eta/\gamma$.
\item Otherwise use the penalty form: covariance $S\trans GS$, expected returns $-\one$,
benchmark $q=S\trans c$, rows $ASz=b$ and $0\le z\le\bar u$. Run the CLA from
$\gamma=\infty$ to $0$. A turning point at $\gamma$ is the knot $\beta=Sz$ at
$\lambda=\gamma$. With $b\ne0$ the start must be a unique vertex
(Section~\ref{sec:limits}).
\item Between two knots $\beta$ is affine in $\lambda$, so any other point on the path is
an interpolation.
\end{enumerate}
In \texttt{cvxcla}, \texttt{Lasso(X, y, a=A)} traces $A\beta=0$ by step~2,
\texttt{nonneg=True} gives $\beta\ge0$, and \texttt{g=G, h=h} gives $G\beta\le h$ with
$h>0$.
\end{minipage}}
\medskip

\begin{figure}[t]
\centering
\includegraphics[width=0.66\textwidth]{portfolio.pdf}
\caption{The diabetes path of Figure~\ref{fig:path}, seen by the budget form as a
portfolio. The long and short positions $Sw$ and the cash weight $w_0$ share a budget
of one, plotted against $t=1/\gamma$. Cash enters at $t=\|\beta(0)\|_1=3460$, the end of
the LASSO path. To its right the portfolio is the least-squares fit diluted by cash.}
\label{fig:portfolio}
\end{figure}

\section{When a reference part gives the zero-sum LASSO}
\label{sec:reference}

The opening example compared the zero-sum LASSO with the log-ratio LASSO against a
reference part $r$. Substituting $\beta_r=-\sum_{j\ne r}\beta_j$ back, the latter is
\[
\min_\beta\ \tfrac12\|y-Z\beta\|_2^2+\lambda\sum_{j\ne r}|\beta_j|
\quad\text{s.t.}\quad\one\trans\beta=0,
\]
the zero-sum LASSO with the penalty on part $r$ removed. The optimality conditions of
the two problems say exactly when they agree. Write $\kappa$ for the multiplier of
$\one\trans\beta=0$ in the zero-sum LASSO and
\[
\rho=Z\trans(y-Z\beta)-\kappa\one ,
\]
so that the zero-sum LASSO's optimality conditions read $\rho=\lambda s$ for a
subgradient $s$ of $\|\beta\|_1$: $s_j=\operatorname{sign}\beta_j$ on the support and
$|s_j|\le1$ off it.

\begin{proposition}[Reference parts]
\label{prop:reference}
Let $\lambda>0$, let $\beta=\beta(\lambda)\ne0$ be the zero-sum LASSO with multiplier
$\kappa$ and $\rho$ as above, and suppose the log-ratio LASSO against part $r$ has a
unique solution at $\lambda$. That solution is $\beta$ if and only if $\beta_r=0$ and
$\rho_r=0$.
\end{proposition}

\begin{proof}
The log-ratio LASSO is optimal at $\beta$ with a multiplier $\kappa'$ exactly when
$Z\trans(y-Z\beta)-\kappa'\one=\lambda s'$ with $s'_j\in\partial|\beta_j|$ for $j\ne r$
and $s'_r=0$, since part $r$ carries no penalty. If $\beta_r=0$ and $\rho_r=0$, take
$\kappa'=\kappa$ and $s'=s$ with $s_r$ replaced by $0$; these satisfy the conditions,
so $\beta$ is optimal, and by uniqueness it is the solution. Conversely, suppose $\beta$
solves the log-ratio LASSO with some $\kappa'$. Then $\lambda s'=\rho+(\kappa-\kappa')\one$.
A nonzero $\beta$ that sums to zero has at least two nonzero coefficients, so one of
them, $\beta_j$ with $j\ne r$, has $s'_j=s_j=\operatorname{sign}\beta_j$, and so
$\kappa'=\kappa$. Then $\rho_r=\lambda s'_r=0$, so $s_r=0$, which rules out
$\beta_r\ne0$.
\end{proof}

On a segment of the zero-sum path the support and signs are fixed, and $\beta$ and
$\kappa$ are affine in $\lambda$, so $\rho_r$ is too. Unless $\rho_r$ vanishes on a
whole segment, it vanishes at most once there. So the log-ratio path against part $r$
meets the zero-sum path at isolated points, at most one per segment, and never where
part $r$ is selected. Off those points the two have different coefficients, though
they can still select the same parts, which is what the matches in
Figures~\ref{fig:example} and~\ref{fig:real} are. The proposition also says where the
points are, so it can be checked: on the opening example $\rho_r$ crosses zero
\exCrossings{} times between knots, and on the real data of Section~\ref{sec:real}
\rdCrossings{} times. At every crossing \texttt{lars\_path} on log-ratios against that
part returns the zero-sum LASSO to \rdCrossAt{} or better; halfway back along the same
segment the two differ by at least \rdCrossOff, relative.

\section{Numerical evidence}
\label{sec:evidence}

Table~\ref{tab:evidence} lists the test problems. The diabetes data \cite{efron2004}
have $n=442$ patients and $p=10$ standardised baseline variables, and the response is
centred. \emph{Nonnegative} requires every coefficient to be at least zero, the one
constraint \texttt{lars\_path} accepts. \emph{Wide} is a Gaussian design with $n=40$ and
$p=100$, with 8 nonzero coefficients and noise of standard deviation $0.5$. \emph{Box}
caps every diabetes coefficient at $|\beta_j|\le u=396.1$, half the largest
least-squares coefficient. At $\lambda=0$ the variables \texttt{bmi} and \texttt{s5}
sit at the cap, and \texttt{s1} does not. \emph{Pinned} requires the fitted value for
the first patient to equal the observed one, one nonhomogeneous equality row. Its
starting vertex puts all weight on \texttt{bmi}, the largest entry of that patient's
covariate vector, and holds still until $\lambda=4053$. \emph{Zero-sum} is the opening
example of Section~\ref{sec:example}, and \emph{Zero-sum, signs} adds nonnegativity on
the even-numbered parts. \emph{sCD14} is the real data of Section~\ref{sec:real}.

Every path is checked in up to four ways. Where LARS applies it is the reference, compared
knot by knot in both $\lambda$ and $\beta$. Elsewhere the reference is an interior-point
solve (Clarabel through CVXPY, tolerances $10^{-12}$) at the midpoint of every segment,
where the path is affine in $\lambda$ and so known exactly from its two ends. This
comparison is only as sharp as the solver, which reaches about $10^{-8}$ relative.
Third, at every knot and every midpoint we evaluate the optimality conditions of
\eqref{eq:lasso} directly, without reference to either form. With
$r=G\beta-c+A\trans\kappa$ they are $r_j=-\lambda\operatorname{sign}\beta_j$ on the
interior of the support, $|r_j|\le\lambda$ at zero, the one-sided versions at a cap, and
primal feasibility. The largest violation, relative to $\|c\|_\infty$, is the last
column of the table. Fourth, where it applies, we compare knot by knot with
\texttt{cvxcla.Lasso} \cite{cvxcla}, a separate implementation: a LASSO homotopy for
sign constraints and rows $G\beta\le h$ with $h>0$, and a rescaled CLA for $A\beta=0$.
It takes six of the eight problems; it has no nonhomogeneous rows, as in
\emph{Pinned}, and no signs on some coefficients only, as in \emph{Zero-sum, signs}.

\begin{table}[t]
\centering
\small
\begin{tabular}{llrrrlccc}
\toprule
problem & form & $n$ & $p$ & knots & reference & gap & \texttt{cvxcla} & KKT \\
\midrule
\input{figures/evidence.tex}
\end{tabular}
\caption{The CLA against LARS (knot by knot) and against an interior-point solver (IPM,
at segment midpoints). \emph{Gap}: largest discrepancy, relative to the largest
coefficient. \emph{\texttt{cvxcla}}: largest knot discrepancy against
\texttt{cvxcla.Lasso} \cite{cvxcla}, where it takes the problem. \emph{KKT}: largest violation of the optimality conditions of
\eqref{eq:lasso} over all knots and midpoints, relative to $\|c\|_\infty$. Each
problem goes through every form that accepts it.}
\label{tab:evidence}
\end{table}

Three observations. First, both forms return the same knots wherever both apply, and
\texttt{cvxcla}'s and LARS's knots to $10^{-13}$, including on the wide design, where 53 knots carry the path
to the interpolating fit. Second, the interior-point gaps are the solver's, not the
path's: on \emph{Box} and \emph{Zero-sum}, \texttt{cvxcla} agrees with the CLA to
$10^{-15}$. The KKT residuals of the CLA's own output are at the level of rounding
throughout, including the cases with no LARS to compare against. Third, the comparison
exposed a quirk of the reference. In positive mode, the last point that
\texttt{lars\_path} returns is labelled $\alpha=0$, but every active correlation there
equals $30.65$. It is the path at $\lambda=30.65$, where the algorithm found nothing
left to enter, rather than the nonnegative least-squares fit at $\lambda=0$. Its KKT
violation at $\lambda=0$ is $3\times10^{-2}$. Both CLA forms carry on to the
nonnegative least-squares solution, and the row in Table~\ref{tab:evidence} compares
that last knot with \texttt{scipy.optimize.nnls} \cite{lawson1974}. All earlier knots
are compared with LARS. Figure~\ref{fig:example}, left, and Figure~\ref{fig:constrained}
show two of the constrained paths that have no LARS reference.

\paragraph{Cost} The exact path is also the cheap option. The baseline is what one
would otherwise do in Python: solve the constrained LASSO with Clarabel through CVXPY at
100 values of $\lambda$, log-spaced over three decades, with the problem compiled once
and $\lambda$ as a parameter, at the solver's default tolerances. That grid costs 30 to
40 times as much as the whole CLA path on \emph{Zero-sum} and \emph{sCD14}, and 170 to
210 times as much on \emph{Box}, and it still misses every knot that falls between two
grid points. On these small problems \texttt{cvxcla.Lasso} costs between $0.7$ and $2$
times what the CLA of this note does. Costs are given as ratios, the fastest of several runs, because
absolute times depend on the machine.

The specialised tool is closer. \texttt{c-lasso} \cite{classo}, run with its path
algorithm on \emph{Zero-sum} and \emph{sCD14}, is a little faster than the CLA, taking
60 to 95 per cent of its time. Its \clPoints{} grid points are close to the path but not
on it: they lie up to \clGap{} from the exact path, relative to the largest coefficient,
where the CLA and \texttt{cvxcla} agree with each other to $10^{-14}$. It takes only
$C\beta=0$, without bounds or signs, and its last release, 1.0.11, needs NumPy~1: it
uses \texttt{np.infty}, which NumPy~2 removed. Path algorithms in other languages, that of
Gaines et al.\ \cite{gaines2018} in MATLAB and \texttt{zeroSum} \cite{zerosum} in R,
are not timed here.

\paragraph{Scale} Table~\ref{tab:scale} runs the budget form on zero-sum problems built
like the opening example, with ten parts in the truth, up to $p=500$ parts and
$n=1000$ samples. An interior-point solve at the midpoint of every segment would cost
more than everything else together, so the checks are the two that scale: the
optimality conditions, and \texttt{cvxcla}. Both hold to rounding at every size. The
100-point grid costs 80 to 110 times as much as the path at every size. The CLA of this
note refactors the bordered system \eqref{eq:border} on every segment, so its cost grows
quickly with the support. The cost of \texttt{cvxcla} grows more slowly: at $p=500$ it
needs a fifth of the time.

\begin{table}[t]
\centering
\small
\begin{tabular}{rrrcc}
\toprule
$n$ & $p$ & knots & \texttt{cvxcla} & KKT \\
\midrule
\input{figures/scale.tex}
\end{tabular}
\caption{The budget form on simulated zero-sum problems of growing size: the largest knot
discrepancy against \texttt{cvxcla.Lasso} and the largest KKT violation over all knots
and midpoints, as in Table~\ref{tab:evidence}.}
\label{tab:scale}
\end{table}

\begin{figure}[t]
\centering
\includegraphics[width=0.66\textwidth]{constrained.pdf}
\caption{The diabetes LASSO under $|\beta_j|\le u$ (dashed), in the penalty form.
Lines: the CLA. Circles: an interior-point solve at the midpoint of every segment. Two
coefficients reach the cap and stay there.}
\label{fig:constrained}
\end{figure}

\section{Real data}
\label{sec:real}

The data come from a study of the gut microbiome in HIV infection
\cite{noguerajulian2016}, as distributed with the R package \texttt{coda4microbiome}
\cite{calle2023}. They are the counts of \rdGenera{} bacterial genera in
\rdSamples{} stool samples, and the response is the level of sCD14, a marker of
inflammation, in thousands. A count of zero has no logarithm, so each count gets a
pseudocount of one, a common choice for such data; the predictors are then the logs of
the shares, centred. The model is the zero-sum LASSO of
Section~\ref{sec:example}.

The budget form traces the whole path in \rdKnots{} knots. Its row in
Table~\ref{tab:evidence}, \emph{sCD14}, shows the same accuracy as on the simulated
problems: the optimality conditions hold to rounding, and \texttt{cvxcla} returns the
same knots. Figure~\ref{fig:real} repeats the experiment of Figure~\ref{fig:example}.
At $\lambda=\rdLambda$ the zero-sum LASSO selects five genera; \rdMatches{} of the
\rdGenera{} reference choices select the same five, and their coefficients are off by up
to \rdWorst\% of the largest one. Nor is this a matter of matching $\lambda$: \rdAnywhere{}
of the \rdGenera{} log-ratio paths select these five genera at any $\lambda$. Over the
first \rdSegments{} segments of the zero-sum path, at most \rdMaxMatches{} reference
choice agrees at the same $\lambda$ and at most \rdMaxAnywhere{} at any $\lambda$, and the
worst coefficient error at the same $\lambda$ never drops below \rdWorstMin\%.

The conclusion does not hang on how the data are prepared. With a pseudocount of one half
instead of one, with $\log$ sCD14 as the response, or with both, at most
\rdSensMaxMatches{} reference choice agrees at the same $\lambda$, and at most
\rdSensMaxAnywhere{} at any $\lambda$, on each of the first \rdSegments{} segments.

\begin{figure}[t]
\centering
\includegraphics[width=\textwidth]{scd14.pdf}
\caption{The zero-sum LASSO of sCD14 on the log-shares of \rdGenera{} genera
\cite{noguerajulian2016}. Left: the path from the CLA in the budget form (lines), with
an interior-point solve at the midpoint of every segment (circles). Right: the
coefficients at $\lambda=\rdLambda$ (dashed line on the left), where the zero-sum LASSO
selects five genera (dots, and ticks). Dashes: \texttt{lars\_path} on log-ratios, one
answer for each of the \rdGenera{} possible reference genera.}
\label{fig:real}
\end{figure}

\section{Where the two forms stop}
\label{sec:limits}

Between them the forms cover coefficient bounds, nonhomogeneous equality rows with a
nondegenerate start, and homogeneous equality rows with sign constraints. The gaps are
all at the start or in the shape of the constraint.

\paragraph{Homogeneous rows in the penalty form} With $b=0$ the starting vertex is
$z=0$, with no free coordinates for $m>0$ rows. This is a degenerate vertex: the
multipliers at the start are not determined by the active set, and \eqref{eq:border}
is singular. The implementation detects this and refuses to start. The budget form is
the remedy, because rescaling to the simplex moves the start off the origin.

\paragraph{Flat starts} A nonhomogeneous row can make the starting linear program
non-unique. The fully invested constraint $\one\trans\beta=1$ of sparse portfolio
regression \cite{brodie2009} is the standard case: every long-only $\beta$ in the
simplex has $\|\beta\|_1=1$ and minimises the $\ell_1$ norm. The path's first segment is
then the least-squares fit over that face, constant in $\lambda$. This is the flat
segment of \cite[Proposition~2]{perspective}, and finding the point requires a
quadratic program rather than a linear one. Markowitz's own treatment of degenerate
starts would handle it, but it is not implemented here.

\paragraph{Inequality rows} Constraints $C\beta\le h$ enter either form as slack assets
with zero variance and zero return, one per row. With $h$ in general position this
changes nothing above. Homogeneous inequalities, such as the ordering constraints of a
monotone fit, make the slack and the bound of the same coordinate active together at
many vertices. This breaks linear independence, and the walk needs the
anti-cycling machinery of the active-set literature. Neither is in the implementation
behind Table~\ref{tab:evidence}, so the note makes no numerical claim for them.

\section{Conclusion}

The LASSO is usually stated, and solved, with no constraint but its $\ell_1$ budget, and
standard path software such as \texttt{lars\_path} cannot take another. The opening
example shows what that costs: the usual workaround for a zero-sum constraint answers a
different question for each choice of reference, and on real microbiome data almost
none of those questions has the zero-sum LASSO's answer. The identity of \cite{perspective}
removes the problem. A constrained LASSO is a portfolio problem, and the Critical Line
Algorithm, in the form Markowitz wrote it, computes its exact path. The penalty form
keeps $\lambda$ and uses the bounds and rows. The budget form keeps the simplex and uses
the scale invariance of homogeneous constraints. On the problems tested, the two forms
agree with each other and with LARS to rounding error, and they meet the LASSO's
optimality conditions to $10^{-15}$ everywhere along the path.

\section*{Code and data availability}

The note's directory holds \texttt{python/make\_figures.py}, a single Python script that
declares its own pinned dependencies, and \texttt{python/make\_figures.py.lock} next to it
fixes Python 3.12 and every package it pulls in, to the exact version and hash. It contains the CLA (the function \texttt{cla}, written for this note rather than taken from
\texttt{cvxcla} \cite{cvxcla}),
the two forms (\texttt{penalty\_form}, \texttt{budget\_form}), the
\texttt{lars\_path} workarounds of Section~\ref{sec:example} and every check, including
the comparison with \texttt{cvxcla}, which the script pins to version 2.3.4.
\texttt{uv run python/make\_figures.py} regenerates the figures,
Table~\ref{tab:evidence} and the numbers of the opening example, and
\texttt{uv run python/make\_figures.py test} runs the checks as tests, and
\texttt{uv run python/make\_figures.py timings} measures the cost ratios.
\texttt{python/classo\_check.py} runs \texttt{c-lasso} 1.0.11 in an environment of its
own, pinned to NumPy~1.26, on the paths \texttt{make\_figures.py} saves, and
\texttt{make figures} runs it after the main script. The diabetes data
ship with scikit-learn, and the simulated problems use fixed seeds. The sCD14 data
are not ours and are not redistributed here: \texttt{python/fetch\_scd14.py} downloads
version 0.2.4 of \texttt{coda4microbiome} from CRAN, checks its SHA-256, and writes the
counts to \texttt{data/scd14.csv}. Every build, in CI as well, runs it first.

\section*{Competing interests}

The author works at Jebel Quant Research, which hosts the code of this note, and
maintains \texttt{cvxcla}, an open-source Python implementation of the Critical Line
Algorithm. The note's results do not depend on it: the CLA it runs is written out in
the companion script, and \texttt{cvxcla} serves only as a second implementation to
check against. The author declares no other competing interests.

\section*{Acknowledgements}

I thank Trevor Hastie for many helpful discussions.

\begin{thebibliography}{99}\setlength{\itemsep}{0pt}

\bibitem{zerosum}
M.~Altenbuchinger, T.~Rehberg, H.~U. Zacharias, et al.
\newblock Reference point insensitive molecular data analysis.
\newblock \emph{Bioinformatics}, 33(2):219--226, 2017.
\newblock \href{https://doi.org/10.1093/bioinformatics/btw598}{\texttt{doi:10.1093/bioinformatics/btw598}}.


\bibitem{brodie2009}
J.~Brodie, I.~Daubechies, C.~De~Mol, D.~Giannone, and I.~Loris.
\newblock Sparse and stable {M}arkowitz portfolios.
\newblock \emph{Proceedings of the National Academy of Sciences}, 106(30):12267--12272,
2009.
\newblock \href{https://doi.org/10.1073/pnas.0904287106}{\texttt{doi:10.1073/pnas.0904287106}}.


\bibitem{calle2023}
M.~L. Calle, M.~Pujolassos, and A.~Sus\'in.
\newblock coda4microbiome: compositional data analysis for microbiome cross-sectional
and longitudinal studies.
\newblock \emph{BMC Bioinformatics}, 24:82, 2023.
\newblock \href{https://doi.org/10.1186/s12859-023-05205-3}{\texttt{doi:10.1186/s12859-023-05205-3}}.


\bibitem{efron2004}
B.~Efron, T.~Hastie, I.~Johnstone, and R.~Tibshirani.
\newblock Least angle regression.
\newblock \emph{The Annals of Statistics}, 32(2):407--499, 2004.
\newblock \href{https://doi.org/10.1214/009053604000000067}{\texttt{doi:10.1214/009053604000000067}}.


\bibitem{gaines2018}
B.~R. Gaines, J.~Kim, and H.~Zhou.
\newblock Algorithms for fitting the constrained lasso.
\newblock \emph{Journal of Computational and Graphical Statistics}, 27(4):861--871, 2018.
\newblock \href{https://doi.org/10.1080/10618600.2018.1473777}{\texttt{doi:10.1080/10618600.2018.1473777}}.


\bibitem{james2020}
G.~M. James, C.~Paulson, and P.~Rusmevichientong.
\newblock Penalized and constrained optimization: an application to high-dimensional
website advertising.
\newblock \emph{Journal of the American Statistical Association}, 115(529):107--122,
2020.
\newblock \href{https://doi.org/10.1080/01621459.2019.1609970}{\texttt{doi:10.1080/01621459.2019.1609970}}.


\bibitem{lawson1974}
C.~L. Lawson and R.~J. Hanson.
\newblock \emph{Solving Least Squares Problems}.
\newblock Prentice-Hall, 1974.

\bibitem{lin2014}
W.~Lin, P.~Shi, R.~Feng, and H.~Li.
\newblock Variable selection in regression with compositional covariates.
\newblock \emph{Biometrika}, 101(4):785--797, 2014.
\newblock \href{https://doi.org/10.1093/biomet/asu031}{\texttt{doi:10.1093/biomet/asu031}}.


\bibitem{markowitz1956}
H.~M. Markowitz.
\newblock The optimization of a quadratic function subject to linear constraints.
\newblock \emph{Naval Research Logistics Quarterly}, 3(1--2):111--133, 1956.
\newblock \href{https://doi.org/10.1002/nav.3800030110}{\texttt{doi:10.1002/nav.3800030110}}.


\bibitem{noguerajulian2016}
M.~Noguera-Julian, M.~Rocafort, Y.~Guill\'en, et al.
\newblock Gut microbiota linked to sexual preference and {HIV} infection.
\newblock \emph{EBioMedicine}, 5:135--146, 2016.
\newblock \href{https://doi.org/10.1016/j.ebiom.2016.01.032}{\texttt{doi:10.1016/j.ebiom.2016.01.032}}.


\bibitem{osborne2000}
M.~R. Osborne, B.~Presnell, and B.~A. Turlach.
\newblock A new approach to variable selection in least squares problems.
\newblock \emph{IMA Journal of Numerical Analysis}, 20(3):389--403, 2000.
\newblock \href{https://doi.org/10.1093/imanum/20.3.389}{\texttt{doi:10.1093/imanum/20.3.389}}.


\bibitem{sklearn}
F.~Pedregosa et al.
\newblock Scikit-learn: machine learning in Python.
\newblock \emph{Journal of Machine Learning Research}, 12:2825--2830, 2011.

\bibitem{perspective}
T.~Schmelzer and T.~Hastie.
\newblock The Critical Line Algorithm and the constrained LASSO: one curve, two
literatures.
\newblock arXiv preprint arXiv:2609.25704, 2026.
\newblock \href{https://doi.org/10.48550/arXiv.2609.25704}{\texttt{doi:10.48550/arXiv.2609.25704}}.


\bibitem{cvxcla}
T.~Schmelzer.
\newblock cvxcla: the Critical Line Algorithm with an operator interface in Python.
\newblock SSRN preprint 7583179, 2026.
\newblock \href{https://doi.org/10.2139/ssrn.7583179}{\texttt{doi:10.2139/ssrn.7583179}}.


\bibitem{lassonote}
T.~Schmelzer.
\newblock The efficient frontier from a LASSO solver.
\newblock arXiv preprint arXiv:2609.37108, 2026.
\newblock \href{https://doi.org/10.48550/arXiv.2609.37108}{\texttt{doi:10.48550/arXiv.2609.37108}}.


\bibitem{classo}
L.~Simpson, P.~L. Combettes, and C.~L. M\"uller.
\newblock c-lasso -- a {P}ython package for constrained sparse and robust regression and
classification.
\newblock \emph{Journal of Open Source Software}, 6(57):2844, 2021.
\newblock \href{https://doi.org/10.21105/joss.02844}{\texttt{doi:10.21105/joss.02844}}.

\bibitem{tibshirani1996}
R.~Tibshirani.
\newblock Regression shrinkage and selection via the lasso.
\newblock \emph{Journal of the Royal Statistical Society, Series B}, 58(1):267--288,
1996.
\newblock \href{https://doi.org/10.1111/j.2517-6161.1996.tb02080.x}{\texttt{doi:10.1111/j.2517-6161.1996.tb02080.x}}.


\end{thebibliography}

\end{document}
